\documentclass[11pt]{article}

% ====================================================
% PACKAGES
% ====================================================
\usepackage[margin=1in, headheight=13.6pt]{geometry}
\usepackage[T1]{fontenc}
\usepackage{lmodern}
\usepackage{microtype}
\usepackage{setspace}
\usepackage{tcolorbox}
\tcbuselibrary{skins,breakable}
\usepackage{amsmath, amssymb, mathtools}
\usepackage{siunitx}

\usepackage{graphicx}
\usepackage{float}
\usepackage{subcaption}

\usepackage{enumitem}
\usepackage{booktabs}

\usepackage{xcolor}
\usepackage{hyperref}

\usepackage{fancyhdr}
\usepackage{lastpage}

\usepackage{listings}
\usepackage{indentfirst}
\usepackage[justification=centering]{caption}

\pagestyle{fancy}
\fancyhf{}
\lhead{\Assignment\ |\ \course}
\rhead{\Name\ |\ \today}
\cfoot{Page \thepage\ of~\pageref{LastPage}}

\hypersetup{
    colorlinks,
    linkcolor={red!50!black},
    citecolor={blue!50!black},
    urlcolor={blue!80!black}
}

\newcommand{\Name}{Arael Anaya}
\newcommand{\Assignment}{Identifying Mission Objectives for a Space Robot}
\newcommand{\course}{Space Robotics}

% ====================================================
% CORRECTION-NEEDED FLAG BOX: loud, red, impossible-to-miss
% marker for text that is now out of date with the settled
% mission design and needs to be rewritten.
% ====================================================
\newtcolorbox{fixme}{
    breakable,
    colback=red!10,
    colframe=red!70!black,
    boxrule=1.2pt,
    arc=2pt,
    left=6pt, right=6pt, top=4pt, bottom=4pt,
    fonttitle=\bfseries\color{red!70!black},
    title={\textcolor{red!70!black}{\Large\textbullet}\ FIX ME --- Out of Date}
}

\begin{document}


% ====================================================
% RESEARCH GAP SELECTION
% ====================================================
\section*{Selected Research Gap}

While most robotic space missions have focused on a single agent carrying all the capabilities needed to complete the
mission objectives, the future of advanced planetary exploration lies in multi-agent systems that are able to share
responsibility and support each other. In order to fully understand a body as complex as Europa, it is necessary to
divide surveying tasks across the main exploration rover and the potential cryobot needed to penetrate the thick layers
of ice. My mission will focus on designing a robust system of agents that can orbit Europa and perform reconnaissance for
a rover traversing the surface in search of a proper drilling site to access the inner layers of the moon. This mission
directly addresses three converging research gaps: autonomous operations in GPS-denied, high-radiation environments
\cite{owlDecadal}; multi-agent coordination and interoperability for cooperative task planning \cite{stmdShortfall};
and fully autonomous approach and proximity operations without Earth-based navigational support \cite{nasaSBIR}.

\subsection*{Mission Objective Statement:}
A coordinated multi-agent system deployed to Europa's surface to autonomously conduct ice characterization and initiate
subsurface access, operating without real-time Earth support.

\begin{fixme}
This statement (and the ``potential cryobot'' framing above) predates the settled mission
design. Every later deliverable (2.3 ConOps, 3.6 Mobility, the 4.1 report) narrows scope to
\emph{identifying} a candidate cryobot insertion site, not ``initiating subsurface access'' or
generating a deployment package --- and settles the agent roster as one rover plus two orbital
relay agents, not a rover-plus-cryobot pairing. Rewrite the Mission Objective Statement to
``...autonomously characterize the ice shell and identify an optimal cryobot insertion site...''
and drop the cryobot-as-agent language in the paragraph above, to match REQ.06.02 and every
downstream document.
\end{fixme}

% ====================================================
% MISSION DESTINATION
% ====================================================
\section*{Mission Destination}

My mission would take place on Europa, Jupiter's icy moon. I chose this destination due to the abundance of knowledge
there is to gain from an in-depth exploration of the system. Europa holds strong indicators of conditions favorable to
life and contains all the chemical building blocks necessary to support it \cite{owlDecadal}. The layers of ice
introduce in-situ resources that might benefit further exploration of the moon as long-term infrastructure is developed,
as well as imposing significant rover design constraints. Any surface vehicle must feature traction-optimized wheels
capable of accurately traversing the fractured, chaotic terrain.

Europa is also incredibly challenging. The 44-minute one-way light delay means every decision from landing onward must
be fully autonomous \cite{nasaSBIR}. Due to the intense radiation environment driven by Jupiter's magnetosphere,
the rover would rely on orbiting satellite robots to scan the terrain and share terrain maps prior to surface traversal.
Such a system would have to be highly robust to faulted or adversarial agents, as incorrect or incomplete maps shared
across the swarm are likely to produce dangerous navigation decisions.

This location is essential to the mission because it provides an incredibly challenging environment that will
necessitate smart, advanced robot design and algorithm development. It is the perfect stage to push robotic
capabilities to their limits.

% ====================================================
% ENVIRONMENTAL & OPERATIONAL CHALLENGES
% ====================================================
\section*{Environmental \& Operational Challenges}

Europa's environment is hostile in ways that shape every design decision in this mission.
Jupiter's magnetosphere floods the moon's surface with high-energy charged particles, exposing
electronics to radiation doses far greater than what commercial components can tolerate. This motivates a requirement
for radiation-hardened electronics and a shielded vault setup for the rover's core computing hardware,
similar to the approach taken on the Galileo spacecraft \cite{owlDecadal}. Redundancy at the agent level
is an equally important mitigation. If one surface relay node builds-up enough single-event upsets to
go silent, the remaining agents must autonomously reconfigure the network to maintain coverage.

Surface temperatures on Europa are around 90~K, which introduces thermal management challenges for both
batteries and actuators. Power storage capacity drops significantly at cryogenic temperatures, meaning
the rover must actively manage its thermal systems to keep cells within an operational range rather
than simply relying on passive insulation.

The complete absence of GPS and the communication blackout due to the 44-minute light delay means
the swarm cannot depend on any external reference for localization or replanning. Each agent must
maintain its own state estimate using onboard sensors and inter-agent ranging, and the swarm must
reach consensus on shared map data using algorithms that are resilient to packet loss and
radiation-induced faults \cite{stmdShortfall, nasaSBIR}. There is no human in the loop to catch
a bad map before the rover commits to a traverse.

% ====================================================
% REQUIRED ROBOT CAPABILITIES
% ====================================================
\section*{Required Robot Capabilities}

\noindent\textbf{Mobility:}

The rover must be able to navigate Europa's icy, fractured terrain and accurately drive without losing
localization. It should be capable of traversing mild inclines and obstacles commonly encountered in
planetary surface missions, with wheel geometry selected for the low-gravity, low-traction conditions
specific to Europa's surface.

\noindent\textbf{Autonomy:}

The robot shall localize itself using onboard sensor data and inter-agent range measurements in the
absence of GPS, generate a planned traverse path, and accurately execute that path. It should also
analyze sensor readings to prioritize target areas based on mission objectives. Ground-penetrating
radar sounder data indicating ice shell thickness will directly inform site selection logic, guiding
the rover toward the thinnest viable ice for cryobot deployment.

\noindent\textbf{Power Systems:}

Due to the rover's distance from the sun, solar power is not viable. Instead, the system must rely on a
Radioisotope Thermoelectric Generator (RTG) or similar nuclear power source. Power budget management
is a crucial. The robot should track consumption across planned traverses and coordinate with the
swarm to avoid scheduling high power draw operations simultaneously across multiple agents.

\noindent\textbf{Communication:}

The robot must communicate with surface relay nodes to exchange sensor readings, map updates, and
mission state. This local mesh network is the primary coordination channel for the swarm. Communication
with Earth still exists as a secondary channel but is non-deterministic due to the 44-minute delay and
periodic occultation by Jupiter. The inter-agent protocol must tolerate partial dropouts and continue
operating under degraded connectivity without stalling consensus.

\noindent\textbf{Durability \& Specialized Instruments:}

The rover shall be designed to withstand Europa's radiation and cryogenic environment, targeting
operation across the full mission lifecycle. It shall carry a ground-penetrating radar sounder (
analogous to the REASON instrument used for Europa Clipper) to characterize ice shell
thickness and build a subsurface map as it traverses the surface in search of the optimal cryobot
insertion point \cite{owlDecadal}.

\newpage

\begin{thebibliography}{9}

\bibitem{owlDecadal}
National Academies of Sciences, Engineering, and Medicine,
\textit{Origins, Worlds, and Life: A Decadal Strategy for Planetary Science and Astrobiology 2023--2032},
The National Academies Press, Washington, D.C., 2022.
[Online]. Available: \url{https://www.nationalacademies.org/read/26522}.
[Accessed: Jun.\ 14, 2026].

\bibitem{stmdShortfall}
NASA Space Technology Mission Directorate,
``Civil Space Technology Shortfall Rankings,''
NASA, Washington, D.C., Jul.\ 2024.
[Online]. Available: \url{https://www.nasa.gov/wp-content/uploads/2024/07/civil-space-shortfall-ranking-july-2024.pdf}.
[Accessed: Jun.\ 14, 2026].

\bibitem{nasaSBIR}
NASA,
``FY 2026--2027 SBIR/STTR Program Solicitation, Appendix 26B-I: SBIR Topics,''
NASA Small Business Innovation Research Program, 2025.

\end{thebibliography}

\end{document}